\section{基本原则}
\begin{enumerate}
  \item 平等原则。考试就考一个问题，平等原则不要求在诉讼中当事人的权利义务相同。
        \begin{enumerate}
          \item 包括权利\zd{相同}。
          \item 也包括权利虽然不同，\zd{相对应}。
          \item \zd{绝对不能狭隘的把平等理解为相同}。
        \end{enumerate}
  \item 同等和对等原则。\zd{只适用于外国}当事人。
  \item 辩论原则。在\zd{诉讼}中，\zd{当事人}有权进行论证。针对的是\zd{事实和证据}。
        \begin{enumerate}
          \item 时间：\zd{贯穿民事诉讼的始终}。
          \item 适用范围：\zd{诉讼程序}。有纠纷你才能辩论，诉讼程序解决纠纷，需要辩论，\zd{其他程序不解决纠纷，不需要辩论}。
          \item 内容：实体和程序都可以。
          \item 形式：口头书面都可以。
        \end{enumerate}
  \item 处分原则。\zd{私权利主体}可以处分自已的私权利，针对的是\zd{诉讼请求层面}。
        \begin{enumerate}
          \item 在法定范围内。
          \item 考试常考的一种违反处分原则的做法，法院的\zd{判决超出原告的诉讼请求}。\qd{多管闲事}。
        \end{enumerate}
        \begin{kk}{约束性辩论原则}
          要让当事人的辩论，对法院的裁判形成约束。
          \begin{enumerate}
            \item 只有当事人提出并\zd{加以阻挡的事实}，法院才能予以认定。
            \item 对当事人都\zd{没有争议的事实}法院应当予以认定 \yjt \zd{受当事人自认约束} \yjt 辩论原则。
            \item 法院对\zd{证据}的调查原则上仅限于当事人提出的 \yjt 辩论原则、处分原则、平等原则。
          \end{enumerate}
        \end{kk}
  \item 诚信原则。\zd{约束所有}参与民事诉讼的主体。
  \item 检查监督原则。\zd{检察机关}作为我国的法律监督机关有权对\zd{民事诉讼活动}进行监督。
        \begin{enumerate}
          \item \zd{监督：把权力关进笼子里}。\zd{对象一定是公权力的行使}。法院 \yjt 审判、执行；
          \item 方式：
                \begin{enumerate}
                  \item 抗诉。只有一种抗诉\zd{再审抗诉}，不满足再审条件不能抗诉。
                  \item 检察建议。再审检察建议。\zd{一定要满足再审条件}；其他检察建议。
                \end{enumerate}
        \end{enumerate}
  \item 在线诉讼原则。谁同意对谁用，不同意各玩各的。
\end{enumerate}

\section{基本制度}
\begin{enumerate}
  \item \zd{合议制}是重点。实践中程序空转，提高司法效率，扩大了独任制的适用范围。
  \item 司法民主。
        \begin{enumerate}
          \item 陪审：事实认定和法律认定分开。
          \item 参审：让陪审员和法官一起组成合议庭，一起对事实的认定、法律的适用问题发表意见，进行评议。
        \end{enumerate}
  \item 审判组织。独任厅\zd{独任制}、合议厅\zd{合议制}、审判委员会。\zd{各种审判组织的适用}：
        \begin{enumerate}
          \item 一审。事实认定。
                \begin{enumerate}
                  \item 简易程序独任制。
                  \item 普通程序合议制。独任制：
                        \begin{itemize}
                          \item 基层法院审理的 + 基本事实清楚、权利关系明确。
                          \item 结论：\zd{一审只要是中院，就不能是独任制}。
                        \end{itemize}
                  \item 人民陪审员只参加适用\zd{一审程序审理的诉讼程序}。
                \end{enumerate}
          \item 二审合议制。吸收当事人对裁判结果不满的功能，事实和法律适用，定纷止争。独任制：
                \begin{enumerate}
                  \item 中院。
                  \item 一审是简易程序 \hz 不服一审裁定。
                  \item 事实清楚、权利关系明确。
                  \item 经双方当事人同意。
                \end{enumerate}
          \item 特别程序独任制，选民资格和重大疑难案件 \yjt 合议制。
          \item 督促程序=支付令 \yjt 独任制。
          \item 公示催告程序。公示催告阶段 \yjt 独任制；除权判决阶段 \yjt 合议制。
        \end{enumerate}
  \item 编口学民诉是一种愚蠢的做法。
  \item 回避。
        \begin{enumerate}
          \item 适用对象：
                \begin{enumerate}
                  \item 与案件有利害关系\zd{（期待保持中立）}。
                  \item \zd{证人，具有不可替代性，虽然需要中立，但他不适用回避制度}。
                \end{enumerate}
          \item 回避的法定原因。很多都是常识问题。
                \begin{enumerate}
                  \item 其他利害关系=能说得出关系，看是否影响公正审理，做题看\zd{有没有修饰性的词语}。
                  \item 有不正当的行为。请客送礼、给当事人推荐律师、违规会见（\zd{非工作时间、工作地点}）当事人。
                  \item 参与过前续审判。例外：发回重审 \yjt 一审 \yjt 二审，\zd{原二审合议庭成员}除外。
                \end{enumerate}
          \item 考试只考申请回避。
                \begin{enumerate}
                  \item 口头书面都可以。
                  \item \zd{开始审理时}提出；开始审理后知道回避事由的，\zd{法庭辩论终结前}。
                  \item 换人\zd{程序不需要重来}；被申请人\zd{暂停工作到决定作出前}，采取紧急措施的除外。
                  \item \zd{审判人员}的回避由院长决定。\zd{判决书有他名字的} \yjt 院长，没名的 \yjt 审判长。
                  \item 提出 \yjt 驳回 \yjt 不服\zd{可以复议一次}，期间\zd{不停止}原决定的执行。
                \end{enumerate}
        \end{enumerate}
  \item 公开审判制度。
        \begin{enumerate}
          \item 法定不公开。国家秘密、个人隐私、其他案件。
          \item 申请不公开。离婚诉讼、商业秘密。一方申请。
          \item \zd{评议一定不能公开}。
          \item \zd{宣判}一定要公开宣判。
        \end{enumerate}
  \item 两审终审制度。
\end{enumerate}
